%!TEX root = ../main_ver2_bench_only.tex
% Appendix

%% ============================================================
\section{Evaluation Dimensions: Detailed Definitions}
\label{app:eval_dimensions}

The \medskillbench evaluation framework scores each skill--task trial along eight binary dimensions.
The first three are designated \emph{core} dimensions; a skill passes overall only if all three core dimensions score~1.
The remaining five are \emph{auxiliary} dimensions that capture execution quality beyond the minimum pass threshold.
Below we provide the full definition and scoring criteria for each dimension.

\paragraph{1.\ \texttt{skill\_invoked} (Core).}
The agent explicitly references, reads, or attempts to use the \skillmd file during the execution trajectory.
A score of~1 is assigned if the agent's trajectory contains evidence that the skill specification was consulted---for example, a file-read action targeting the \skillmd, direct quotation of skill instructions, or an explicit statement acknowledging the skill.
A score of~0 indicates the agent completed the task without any evidence of engaging with the provided skill.

\paragraph{2.\ \texttt{skill\_used\_correctly} (Core).}
The agent applies the skill in a manner consistent with the specification's intended use case, input--output contract, and procedural instructions.
A score of~1 requires that the agent's actions align with the skill's documented workflow---correct parameter usage, appropriate invocation context, and adherence to any constraints or prerequisites stated in the specification.
Partial or superficial references to the skill that do not translate into specification-aligned actions receive a score of~0.

\paragraph{3.\ \texttt{task\_completion} (Core).}
The agent produces a final output that satisfies the task objective.
Scoring considers whether the output addresses the stated goal, provides a substantive answer or artifact, and is not truncated, empty, or obviously incorrect.
A score of~1 does not require perfect accuracy but demands that the output constitutes a reasonable, complete attempt.

\paragraph{4.\ \texttt{tool\_execution\_success} (Auxiliary).}
At least one tool invocation during the trajectory produces a verified, non-trivial output.
This dimension captures whether the agent successfully grounds the skill specification into executable tool calls that return meaningful results.
A score of~0 indicates that all tool calls either failed, produced empty outputs, or were never attempted despite the skill specifying tool-based workflows.

\paragraph{5.\ \texttt{output\_quality} (Auxiliary).}
The final output demonstrates domain-appropriate quality: accurate terminology, logically coherent reasoning, and sufficient depth for the task's complexity.
The judge assesses whether the output would be considered adequate by a domain practitioner, accounting for the difficulty of the generated task.

\paragraph{6.\ \texttt{reasoning\_quality} (Auxiliary).}
The agent's reasoning trajectory---including intermediate steps, tool-call rationale, and error recovery---demonstrates logical coherence and appropriate problem decomposition.
A score of~1 indicates that the agent's chain of thought is traceable, non-circular, and reflects genuine engagement with the problem rather than superficial pattern matching.

\paragraph{7.\ \texttt{efficiency} (Auxiliary).}
The agent completes the task within a reasonable number of iterations relative to task complexity, without excessive redundant actions, repeated failures on the same step, or unnecessary detours.
A score of~0 is assigned when the agent exhausts the iteration budget, enters retry loops without progress, or performs substantially more actions than necessary.

\paragraph{8.\ \texttt{trajectory\_quality} (Auxiliary).}
The overall trajectory exhibits purposeful progression toward the goal, appropriate error handling, and coherent state management across iterations.
This holistic dimension captures trajectory-level phenomena not addressed by individual dimensions: does the agent recover from failures constructively, maintain context across steps, and converge toward a solution?

\paragraph{Aggregation.}
A skill achieves \emph{pass} status when all three core dimensions score~1.
A skill achieves \emph{gold} status when all eight dimensions score~1.
\emph{Silver} status denotes core pass with at least one auxiliary failure.
The \rbu gap is computed as the difference between \texttt{skill\_invoked} and \texttt{task\_completion} (see Equation~\ref{eq:rbu} in the main text).


%% ============================================================
\section{Mandatory Specification-Access Instruction}
\label{app:read_first_prompt}

Every skill evaluation trial injects the following system-level instruction before the task description to ensure the agent engages with the provided skill specification.
This prompt is appended to the agent's system message and is not visible to the judge.

\begin{small}
\begin{verbatim}
IMPORTANT: Before attempting this task, you MUST
consult the provided skill specification document.
This document contains critical instructions,
tool specifications, and domain knowledge required
to complete the task correctly.

Steps:
1. Review the full skill specification before taking any action.
2. Follow the procedures, constraints, and tool
   invocations described in the specification.
3. If the specification identifies particular tools, libraries,
   or commands, use them as directed.
4. Only after reviewing and understanding the specification
   should you proceed to address the task.

Do NOT skip the specification. Do NOT rely solely on
your internal knowledge if the specification provides
specific procedures.
\end{verbatim}
\end{small}

This prompt achieves a 96.2\% invocation rate across the full corpus of 1,462 skills, validating its effectiveness at directing agent attention.
The remaining 3.8\% of non-invocation cases arise primarily from skills whose specification documents are empty, malformed, or contain only metadata headers without actionable content.


%% ============================================================
\section{Skill Domain Distribution}
\label{app:domain_distribution}

The 917 Tier-A skills in \medskillbench span multiple medical and biomedical domains, reflecting the breadth of open-source medical agent tool repositories.
Domain annotations are produced by an automated classifier with manual verification of borderline cases.

\paragraph{Primary domains.}
The corpus covers the following major areas:
\begin{itemize}[nosep,leftmargin=*]
  \item \textbf{Biology \& Bioinformatics} --- Sequence analysis (BLAST, multiple sequence alignment), gene expression analysis, phylogenetics, genome annotation, protein structure prediction, and pathway analysis tools.
  \item \textbf{Clinical Medicine} --- Clinical decision support, diagnostic reasoning guides, treatment protocol references, drug interaction checkers, and clinical calculator wrappers.
  \item \textbf{Drug Development \& Pharmacology} --- Molecular docking, ADMET prediction, compound screening, pharmacokinetic modeling, and drug--target interaction databases.
  \item \textbf{Chemoinformatics} --- Chemical structure manipulation (RDKit-based), molecular fingerprinting, SMILES/InChI conversion, and chemical property prediction.
  \item \textbf{Proteomics \& Structural Biology} --- Protein--protein interaction analysis, mass spectrometry data processing, structural alignment, and binding site prediction.
  \item \textbf{Genomics \& Variant Analysis} --- Variant calling, VCF processing, genome-wide association analysis, and annotation pipelines.
  \item \textbf{Medical Literature \& Evidence} --- PubMed search, systematic review automation, evidence grading, and citation management.
  \item \textbf{Laboratory \& Safety} --- Laboratory protocol guides, biosafety references, equipment operation procedures, and regulatory compliance checklists.
  \item \textbf{Medical Imaging} --- DICOM processing, radiological analysis aids, and image segmentation tool wrappers.
  \item \textbf{Public Health \& Epidemiology} --- Epidemiological modeling, outbreak analysis, surveillance data processing, and population health statistics.
\end{itemize}

\paragraph{Distribution characteristics.}
The domain distribution is intentionally unbalanced, reflecting the natural distribution of skills in the source repositories.
Biology \& Bioinformatics and Clinical Medicine together account for the largest share, followed by Drug Development and Medical Literature tools.
This imbalance is preserved rather than artificially balanced, as it reflects the real-world landscape of available medical agent tools.


%% ============================================================
\section{Full Benchmark Statistics}
\label{app:full_benchmark}

\Cref{tab:full_benchmark} extends the summary in Table~\ref{tab:benchmark_match} with additional columns for absolute match counts, unique skill counts, and the top match types observed per benchmark.

\begin{table*}[t]
\centering\small
\caption{Extended skill--benchmark matching statistics across 12 medical QA benchmarks (917 Tier-A skills).
``Matched'' indicates the number of questions with at least one skill match.
``Top match types'' lists the two most frequent match type annotations for each benchmark.}
\begin{tabular}{lrrrrll}
\toprule
\textbf{Benchmark} & \textbf{Total Q} & \textbf{Matched} & \textbf{Match\%} & \textbf{Skills} & \textbf{Top Match Type 1} & \textbf{Top Match Type 2} \\
\midrule
Lab-Bench         &   900 & 416 & 46.2 & 90 & \texttt{database\_query}    & \texttt{literature\_search} \\
MedAgents         &   332 &  96 & 28.9 & 16 & \texttt{literature\_search}  & \texttt{reasoning\_guide}   \\
MedMCQA           & 4,183 &1,071& 25.6 & 50 & \texttt{literature\_search}  & \texttt{database\_query}    \\
PubMedQA          & 1,000 & 194 & 19.4 & 35 & \texttt{literature\_search}  & \texttt{database\_query}    \\
HeadQA            & 2,742 & 499 & 18.2 & 66 & \texttt{literature\_search}  & \texttt{reasoning\_guide}   \\
SGI-Reasoning     &    72 &   9 & 12.5 & 17 & \texttt{reasoning\_guide}    & \texttt{literature\_search} \\
MMLU-Med          & 1,089 & 127 & 11.7 & 42 & \texttt{literature\_search}  & \texttt{reasoning\_guide}   \\
MedQA-USMLE       & 1,273 &  71 &  5.6 & 21 & \texttt{literature\_search}  & \texttt{reasoning\_guide}   \\
LabSafety         &   765 &  42 &  5.5 &  3 & \texttt{literature\_search}  & \texttt{reasoning\_guide}   \\
MedCalc-Bench     & 1,067 &  47 &  4.4 &  5 & \texttt{database\_query}     & \texttt{reasoning\_guide}   \\
MedBullets-5op    &   298 &  15 &  5.0 &  6 & \texttt{literature\_search}  & \texttt{reasoning\_guide}   \\
MedBullets-4op    &   298 &   8 &  2.7 &  6 & \texttt{literature\_search}  & \texttt{reasoning\_guide}   \\
\midrule
\textbf{Overall}  & \textbf{14,019} & \textbf{2,596} & \textbf{18.5} & \textbf{178} & \texttt{literature\_search} (71.7\%) & \texttt{database\_query} (21.0\%) \\
\bottomrule
\end{tabular}
\label{tab:full_benchmark}
\end{table*}

\paragraph{Notes on matched counts.}
The ``Matched'' column reports unique questions with at least one skill assignment; a single question may be matched to multiple skills, yielding approximately 3,898 total skill--question pairs across the 2,596 matched questions.
Relevance annotations indicate that 91.5\% of matches are classified as \emph{weak} relevance (the skill provides tangential support) and 8.5\% as \emph{strong} relevance (the skill directly addresses the core reasoning required by the question).

\paragraph{Skill concentration analysis.}
Across all benchmarks, the top three matched skills are \texttt{agent-browser} (appearing in 78\% of all matched pairs), \texttt{ai-analyzer} (13.8\%), and \texttt{bio-entrez-search} (12.2\%).
This heavy concentration underscores that current medical skill repositories favor general-purpose retrieval and analysis tools over domain-specialized reasoning aids.
Lab-Bench is an exception, with 90 unique skills matched---the highest skill diversity of any benchmark---consistent with its emphasis on specific laboratory protocols and bioinformatics pipelines.


%% ============================================================
\section{Evaluation Stability}
\label{app:eval_stability}

To assess evaluation stability across the corpus, we report results stratified by processing batch in \Cref{tab:batch_stability}.
Skills were processed in sequential batches of 400 (with the final batch containing the remaining 262 skills).

\begin{table}[t]
\centering\small
\caption{Evaluation results by processing batch, demonstrating stability across the corpus.}
\begin{tabular}{lrccr}
\toprule
\textbf{Batch} & \textbf{$n$} & \textbf{Pass\%} & \textbf{Gold\%} & \textbf{\rbu Gap} \\
\midrule
1 (1--400)       & 400 & 57.8 & 15.0 & 0.333 \\
2 (401--800)     & 400 & 65.5 & 23.8 & 0.173 \\
3 (801--1200)    & 400 & 68.3 & 22.8 & 0.168 \\
4 (1201--1462)   & 262 & 67.2 & 20.2 & 0.199 \\
\midrule
\textbf{Overall} & \textbf{1,462} & \textbf{64.4} & \textbf{20.5} & \textbf{0.220} \\
\bottomrule
\end{tabular}
\label{tab:batch_stability}
\end{table}

Batches 2--4 show stable performance, with pass rates in the 65.5--68.3\% range and \rbu gaps between 0.168 and 0.199, confirming that the evaluation framework yields consistent results across the corpus.


%% ============================================================
\section{Skill-Augmented QA: Detailed Results}
\label{app:ablation_details}

This appendix provides extended results for the skill-augmented QA evaluation reported in \Cref{ssec:ablation_results}.

\paragraph{Per-benchmark breakdown.}
\Cref{tab:ablation_full} reports per-benchmark accuracy under both conditions for all evaluated models, including benchmarks omitted from the main-text table due to small matched sample sizes.

\begin{table*}[t]
\centering\small
\caption{Extended skill-augmented QA results across all 12 medical benchmarks. $n$: number of evaluated (question, skill) pairs per benchmark. Bold $\Delta$ indicates $p < 0.05$ (McNemar's test).}
\begin{tabular}{lrcccccc}
\toprule
\textbf{Benchmark} & \textbf{$n$} & \multicolumn{3}{c}{\textbf{GPT-5.1}} & \multicolumn{3}{c}{\textbf{Claude Sonnet~4.5}} \\
\cmidrule(lr){3-5} \cmidrule(lr){6-8}
 & & w/o & w/ & $\Delta$ & w/o & w/ & $\Delta$ \\
\midrule
Lab-Bench       & \todo{} & \todo{} & \todo{} & \todo{} & \todo{} & \todo{} & \todo{} \\
MedAgents       & \todo{} & \todo{} & \todo{} & \todo{} & \todo{} & \todo{} & \todo{} \\
MedMCQA         & \todo{} & \todo{} & \todo{} & \todo{} & \todo{} & \todo{} & \todo{} \\
PubMedQA        & \todo{} & \todo{} & \todo{} & \todo{} & \todo{} & \todo{} & \todo{} \\
HeadQA          & \todo{} & \todo{} & \todo{} & \todo{} & \todo{} & \todo{} & \todo{} \\
SGI-Reasoning   & \todo{} & \todo{} & \todo{} & \todo{} & \todo{} & \todo{} & \todo{} \\
MMLU-Med        & \todo{} & \todo{} & \todo{} & \todo{} & \todo{} & \todo{} & \todo{} \\
MedQA-USMLE     & \todo{} & \todo{} & \todo{} & \todo{} & \todo{} & \todo{} & \todo{} \\
LabSafety       & \todo{} & \todo{} & \todo{} & \todo{} & \todo{} & \todo{} & \todo{} \\
MedCalc-Bench   & \todo{} & \todo{} & \todo{} & \todo{} & \todo{} & \todo{} & \todo{} \\
MedBullets-5op  & \todo{} & \todo{} & \todo{} & \todo{} & \todo{} & \todo{} & \todo{} \\
MedBullets-4op  & \todo{} & \todo{} & \todo{} & \todo{} & \todo{} & \todo{} & \todo{} \\
\midrule
\textbf{Overall} & \todo{} & \todo{} & \todo{} & \todo{} & \todo{} & \todo{} & \todo{} \\
\bottomrule
\end{tabular}
\label{tab:ablation_full}
\end{table*}

\paragraph{Statistical details.}
For each benchmark--model pair, we report the McNemar $\chi^2$ statistic and two-sided $p$-value.
Bootstrap 95\% confidence intervals for the uplift are computed via 10,000 resamples of the paired (with, without) binary outcome vectors.
\todo{Populate statistical results once Step~2 experiments complete.}

\paragraph{Uplift distribution.}
Per-question uplift takes one of three values: $+1$ (with-skill correct, without-skill incorrect), $0$ (same outcome), or $-1$ (with-skill incorrect, without-skill correct).
The proportion of $+1$ vs.\ $-1$ outcomes determines the direction and magnitude of the aggregate effect.
\todo{Report uplift distribution breakdown and add histogram figure.}


%% ============================================================
\section{SGI Sensitivity Analysis}
\label{app:sgi_sensitivity}

We verify that the core findings are robust to the choice of $w_c$ in the \textsc{SGI} definition (\Cref{eq:sgi}).
The dominance principle (\Cref{eq:wc_derivation}) requires $w_c > 0.75$ to guarantee that any core-failing skill ranks below any core-passing skill.
We report results for $w_c \in \{0.60, 0.70, 0.75, 0.80\}$.

\begin{table}[t]
\centering\small
\caption{\textsc{SGI} sensitivity to $w_c$. Domain rankings and the separation between core-passing and core-failing skills are reported for each weight.}
\begin{tabular}{lcccc}
\toprule
\textbf{Metric} & $w_c{=}0.60$ & $w_c{=}0.70$ & $w_c{=}0.75$ & $w_c{=}0.80$ \\
\midrule
Mean \textsc{SGI} (all)          & \todo{} & \todo{} & \todo{} & \todo{} \\
Mean \textsc{SGI} (core-pass)    & \todo{} & \todo{} & \todo{} & \todo{} \\
Mean \textsc{SGI} (core-fail)    & \todo{} & \todo{} & \todo{} & \todo{} \\
Separation (pass $-$ fail)       & \todo{} & \todo{} & \todo{} & \todo{} \\
Domain rank correlation ($\tau$) & \todo{} & \todo{} & \todo{} & \todo{} \\
\bottomrule
\end{tabular}
\label{tab:sgi_sensitivity}
\end{table}

\todo{Compute SGI under all four $w_c$ values from Step~0 data and verify that domain rankings are stable (Kendall's $\tau > 0.9$).}
At $w_c = 0.60$, the dominance property is violated: a core-failing skill with perfect auxiliary quality achieves $\textsc{SGI} = 0.80$, outscoring a core-passing skill with zero auxiliary quality ($\textsc{SGI} = 0.60$).
At $w_c \geq 0.75$, the separation is guaranteed.


%% ============================================================
\section{Statistical Testing Details}
\label{app:statistical_details}

This appendix provides the full statistical testing framework for the skill-augmented QA evaluation.

\paragraph{McNemar's test for paired binary outcomes.}
For each skill--question pair $(x_i, s, g_i) \in \mathcal{P}_\mu$, let $c_{i,s}^{+} = \mathbf{1}[\mathcal{M}(x_i, s) = g_i]$ (with-skill correctness) and $c_{i,s}^{-} = \mathbf{1}[\mathcal{M}(x_i, \varnothing) = g_i]$ (without-skill correctness).
The discordant counts are:
\begin{equation}
\begin{aligned}
b &= \sum_{(i,s)} \mathbf{1}[c_{i,s}^{+} = 1 \wedge c_{i,s}^{-} = 0], \\
c &= \sum_{(i,s)} \mathbf{1}[c_{i,s}^{+} = 0 \wedge c_{i,s}^{-} = 1].
\end{aligned}
\label{eq:mcnemar_counts}
\end{equation}
Under $H_0$: skill availability has no effect, $b$ and $c$ follow a binomial distribution with equal probabilities.
The McNemar statistic is:
\begin{equation}
\chi^2_{\text{McN}} = \frac{(b - c)^2}{b + c} \;\sim\; \chi^2(1).
\label{eq:mcnemar_stat}
\end{equation}
We reject $H_0$ at $\alpha = 0.05$ if $\chi^2_{\text{McN}} > 3.841$.

\paragraph{Multi-match dependence.}
When a question matches multiple skills ($|\mu(x_i)| > 1$), pairs sharing the same question have correlated without-skill outcomes $c_{i,s}^{-}$, violating the independence assumption.
We address this via clustered bootstrap (clustering by question index $i$) as a robustness check alongside the standard McNemar test.

\paragraph{Bootstrap confidence intervals.}
For $\Delta_{\text{uplift}}$, we construct bootstrap 95\% CIs by resampling the paired outcome vectors with replacement $B = 10{,}000$ times, clustering by question index $i$:
\begin{equation}
\text{CI}_{95\%}(\Delta_{\text{uplift}}) = \left[\hat{\Delta}^{*(0.025)}, \;\hat{\Delta}^{*(0.975)}\right],
\label{eq:bootstrap_ci}
\end{equation}
where $\hat{\Delta}^{*(\alpha)}$ is the $\alpha$-quantile of the bootstrap distribution.

\paragraph{Groundability-uplift correlation.}
The aggregate interface-type correlation (\Cref{eq:uplift_prediction}) is computed over $|\mathcal{I}| = 5$ data points.
With $n = 5$, the critical value for significance at $\alpha = 0.05$ is $|r| > 0.878$, yielding negligible statistical power.
We therefore use the per-skill \textsc{SGI}$\to$uplift regression (\Cref{eq:per_skill_reg}) as the primary test, which operates over individual skills (higher $n$) and provides substantially greater power.
\todo{Report both $r_{GU}$ and $\hat{\beta}_1$ with standard errors once Step~2 experiments complete.}
